\section{Akar Filosofis: Dari Silogisme hingga Logika Simbolik}

Sebelum era di mana mesin silikon mampu memproses triliunan data per detik, fondasi dari ilmu komputer (*Computer Science*) sesungguhnya diletakkan pertama kali bukan oleh para insinyur kabel tembaga, melainkan oleh para filsuf kuno yang bermimpi merumuskan \textbf{"Hukum Pasti Berpikir"} \cite{wiki_proplogic}. 

\logic{Logika Matematika} sejatinya adalah instrumen purba umat manusia untuk merebus bahasa alamiah yang penuh emosi dan ambiguitas menjadi sebuah rumus aljabar yang mutlak, tak terbantahkan, dan siap dieksekusi oleh mesin buta sekalipun \cite{rosen2019,copi2014,forallx}.

\subsection{1. Fajar Logika Yunani: Era Silogisme Aristoteles}

Pada rentang tahun 384--322 SM, Aristoteles merintis mahakarya intelektual pertama yang membakukan struktur perdebatan manusia menjadi rumus deduksi kaku yang disebut \textbf{Silogisme} \cite{iep_logic_history}. Era ini menandai terobosan di mana manusia menyadari kebenaran suatu opini bisa dinilai sah hanya dengan melacak struktur tata bahasanya, terlepas dari omong kosong apa isi aslinya.

\begin{contoh}[Bentuk Klasik Silogisme Aristoteles]
\indent \textbf{Premis Mayor:} Semua mamalia wajib mengembuskan karbon dioksida. \\
\indent \textbf{Premis Minor:} Kucing perak saya tergolong bangsa mamalia. \\
\indent \textbf{Validasi Kesimpulan:} Oleh mutlak karenanya, Kucing perak saya pastilah mengembuskan karbon dioksida.
\end{contoh}

Sumbangan kaum Yunani Kuno (khususnya penerus filosofi Stoik) melahirkan dasar-dasar penghubung kondisional "Jika-Maka" (\textit{If-Then}). Sayangnya, selama ribuan tahun, analisis logika mereka masih terbelenggu di dalam kurungan huruf-huruf abjad lisan manusia, belum sepenuhnya matematis transparan.

\subsection{2. Era Revolusi Simbolik: George Boole \& Gottlob Frege}

Pada abad ke-19, logika resmi memisahkan diri dari belenggu sastra perdebatan sosial dan dikawinkan paksa dengan ketegasan operasi Matematika \cite{openlogic,copi2014}. Aktor terbesarnya adalah matematikawan eksentrik Inggris, \textbf{George Boole (1815-1864)}. 

Dalam mahakaryanya yang menggetarkan dunia, \textit{The Laws of Thought (1854)}, Boole secara revolusioner mendeklarasikan gagasan gila: *"Bagaimana jika nilai kebenaran suatu kalimat direbus ke dalam rumus hitung matematika, persis layaknya operasi Tambah (+) dan Kali (x)?"* Boole menciptakan cikal bakal \textbf{Aljabar Boolean}. Dialah Bapak yang memformat dunia alam semesta bahasa menjadi mutlak diukur cuma dalam skor biner **1 (True) dan 0 (False)**.

Estafet kejeniusan ini disambut oleh \textbf{Gottlob Frege (1879)} melalui karyanya \textit{Begriffsschrift} \cite{copi2014}. Frege menyeret logika Boole menyentuh dewa langit dengan menciptakan amunisi \textbf{Logika Predikat}. Ia memperkenalkan panah Kuantifier Universal "Semua" ($\forall$) dan Eksistensial "Sebagian" ($\exists$). Berkat beliau, algoritma Database Relasional jutaan triliun baris (seperti SQL) bisa melahap pencarian kondisi majemuk dengan akurasi 100\% tak meleset se-bit pun di industri digital masa ini.

\begin{figure}[htbp]
\centering
\includegraphics[width=0.9\textwidth]{timeline-logika}
\caption{Timeline evolusi Logika Matematika: Dari debu perdebatan filsuf, menembus ke ranah silikon transistor cerdas \cite{wiki_proplogic}.}
\label{fig:timeline-logika}
\end{figure}
